% Programming with InitechBase, Chapter 7: Working with the Database (AN-16;
% annex Sec 8.2 P10/P11).  Panels replayed (examples/p7-*.dot); sources
% cmd/query.c (SEEK, FOUND), cmd/mutate.c (APPEND BLANK, REPLACE, index
% maintenance), samir_main.c repl_do_use (INDEX).
\chapter{Working with the Database}
\Gates{test-pubs-examples, test-interp-replace, test-samir-query, test-ndx-maintain}
The programs so far read the ledger.  This chapter adds invoices to it: a
program asks for the invoice, makes sure the number is new, and adds the
record.

\begin{ChapterCovers}
\item The databases a program works on
\item Finding a record with \cmd{SEEK}
\item Adding a record: \cmd{APPEND BLANK} and \cmd{REPLACE}
\item Index files in programs
\item Entering an invoice
\end{ChapterCovers}

\ATEHead{Preparing for This Chapter}
Start InitechBase with a fresh copy of the sample diskette.

\ATEHead{The Databases of a Program}
A program works on the databases opened at the dot prompt before it runs.  The
programs of the system expect \file{INVOICE.DBF} in the selected work area; the
programs that find invoices by number also expect its index \file{INVNO.NDX}.
Each program says so in its comment lines.  A program can select other work
areas with \cmd{SELECT}.

\ATEHead{Index Files in Programs}
The index file named first in \cmd{USE} gives the order in which a program
meets the records, and is the index \cmd{SEEK} searches.  The ledger has two
indexes, on the invoice number and on the customer:
\DotExample{examples/p7-order}
The index files open with a database are kept up to date when a program adds
or changes records.

\ATEHead{Entering an Invoice}
\file{NEWINV.PRG} asks for the invoice number, looks it up, and stops if the
number is already in the ledger.  Otherwise it asks for the rest of the invoice,
adds a blank record and fills it in:
\ProgramListing{NEWINV.PRG}
\begin{Steps}
\item Type \type{USE INVOICE INDEX INVNO}.
\item Type \type{DO NEWINV}.
\item Answer the questions: \type{B1005}, \type{Initrode}, \type{420.00},
\type{04/15/92}.
\item Type \type{LIST INVNO, CUST, AMOUNT, DUEDATE}.
\end{Steps}
\DotExample{examples/p7-newinv}
A number already in the ledger is refused:
\DotExample{examples/p7-dup}

\ATEHead{How the Program Works}
\SideHead{Checking first}The program seeks the new number before it adds
anything.  \cmd{RETURN} inside the \cmd{IF} ends the program at once when the
number is found.

\SideHead{The answers}\cmd{ACCEPT} takes the customer and the due date as
characters, \cmd{INPUT} the amount as a number.

\SideHead{The record}\cmd{APPEND BLANK} adds the record and makes it current;
two \cmd{REPLACE} commands fill in its fields.  \fn{UPPER()} keeps the customer in
capitals and \fn{CTOD()} converts the due date.  The record is written to the
disk at once, and the index is brought up to date.

\ATEHead{Where Do You Go From Here?}
Chapter~8 prints the ledger as a report.

\ATEHead{Summary}
\begin{enumerate}
\item Open the databases and indexes a program needs before it runs.
\item \cmd{SEEK} and \fn{FOUND()} check whether a key is present.
\item \cmd{APPEND BLANK} and \cmd{REPLACE} add a record; the open indexes are
  kept up to date.
\end{enumerate}
